\documentclass[11pt]{article}
\usepackage[utf8]{inputenc}
\usepackage[spanish, es-tabla]{babel}
\usepackage{amsmath, amssymb}
\usepackage{booktabs}
\usepackage{enumitem}
\usepackage{geometry}
\geometry{letterpaper, margin=2.5cm}

\title{Econometría ICS-294\\Guía de ejercicios: Regresión Lineal Simple}
\author{Francisco Alfaro Medina}
\date{}

\begin{document}
\maketitle

\section*{Instrucciones}
Responda de forma clara, mostrando fórmulas, cálculos e interpretación económica. Cuando use MCO, indique explícitamente cuál es la variable dependiente y cuál es la variable explicativa.

\section*{Ejercicios}

\begin{enumerate}[label=\textbf{Ejercicio \arabic*.}, leftmargin=*]

    \item \textbf{Modelo de regresión lineal simple.}
    
    Considere el modelo poblacional:
    \[
        y=\beta_0+\beta_1x+u.
    \]
    \begin{enumerate}[label=\alph*)]
        \item Explique qué representan $y$, $x$, $\beta_0$, $\beta_1$ y $u$.
        \item ¿Qué significa interpretar $\beta_1$ ``ceteris paribus''?
        \item Proponga un ejemplo económico donde $y$ sea salario y $x$ sea educación. Mencione al menos tres variables que podrían estar dentro de $u$.
    \end{enumerate}

    \item \textbf{Esperanza condicional y FRP.}
    
    En una población pequeña se observa la siguiente relación entre años de educación e ingreso:
    
    \begin{center}
    \begin{tabular}{lcc}
    \toprule
    Persona & Educación & Ingreso \\
    \midrule
    A & 10 & 400 \\
    B & 10 & 500 \\
    C & 12 & 600 \\
    D & 12 & 700 \\
    E & 14 & 900 \\
    F & 14 & 1000 \\
    \bottomrule
    \end{tabular}
    \end{center}
    
    \begin{enumerate}[label=\alph*)]
        \item Calcule $E[Y|X=10]$, $E[Y|X=12]$ y $E[Y|X=14]$.
        \item Dibuje intuitivamente la función de esperanza condicional.
        \item Si la FRP fuera lineal, ¿qué relación tendría con $E[Y|X]$?
    \end{enumerate}

    \item \textbf{Estimación manual por MCO.}
    
    Use los siguientes datos:
    
    \begin{center}
    \begin{tabular}{cccccc}
    \toprule
    $i$ & 1 & 2 & 3 & 4 & 5 \\
    \midrule
    $x_i$ & 1 & 2 & 3 & 4 & 5 \\
    $y_i$ & 2 & 3 & 5 & 4 & 6 \\
    \bottomrule
    \end{tabular}
    \end{center}
    
    \begin{enumerate}[label=\alph*)]
        \item Calcule $\bar{x}$ y $\bar{y}$.
        \item Calcule $\hat{\beta}_1=\frac{\sum_i(x_i-\bar{x})(y_i-\bar{y})}{\sum_i(x_i-\bar{x})^2}$.
        \item Calcule $\hat{\beta}_0=\bar{y}-\hat{\beta}_1\bar{x}$.
        \item Escriba la función de regresión muestral.
        \item Interprete $\hat{\beta}_1$.
    \end{enumerate}

    \item \textbf{Valores ajustados y residuos.}
    
    Usando la regresión estimada en el ejercicio anterior:
    \begin{enumerate}[label=\alph*)]
        \item Calcule $\hat{y}_i$ para cada observación.
        \item Calcule el residuo $\hat{u}_i=y_i-\hat{y}_i$.
        \item Verifique que $\sum_i \hat{u}_i=0$.
        \item Verifique que $\sum_i x_i\hat{u}_i=0$.
    \end{enumerate}

    \item \textbf{Propiedades aritméticas de MCO.}
    
    Explique por qué, en una regresión con constante:
    \begin{enumerate}[label=\alph*)]
        \item la recta estimada pasa por el punto $(\bar{x},\bar{y})$;
        \item el promedio de los valores ajustados es igual al promedio de $y$;
        \item los residuos tienen promedio cero.
    \end{enumerate}

    \item \textbf{Bondad de ajuste.}
    
    Para una regresión se obtiene:
    \[
        SST=500,\qquad SSR=125.
    \]
    \begin{enumerate}[label=\alph*)]
        \item Calcule $SSE$.
        \item Calcule $R^2$ usando $R^2=SSE/SST$.
        \item Calcule $R^2$ usando $R^2=1-SSR/SST$.
        \item Interprete el resultado.
    \end{enumerate}

    \item \textbf{Interpretación de una regresión estimada.}
    
    Suponga que se estima la relación entre salario por hora y años de educación:
    \[
        \widehat{salario}=0.095+0.54\,educ.
    \]
    \begin{enumerate}[label=\alph*)]
        \item Interprete el intercepto.
        \item Interprete la pendiente.
        \item Prediga el salario para una persona con 12 años de educación.
        \item ¿Qué limitaciones tiene interpretar esta pendiente como causal?
    \end{enumerate}

    \item \textbf{Media condicional cero.}
    
    Considere nuevamente:
    \[
        salario=\beta_0+\beta_1educ+u.
    \]
    \begin{enumerate}[label=\alph*)]
        \item Explique qué significa $E[u|educ]=0$.
        \item Mencione dos variables omitidas que podrían hacer que este supuesto falle.
        \item Explique intuitivamente por qué, si $educ$ está correlacionada con factores dentro de $u$, MCO puede entregar una relación espuria.
    \end{enumerate}

    \item \textbf{Supuestos de RLS.}
    
    Para cada caso, indique qué supuesto podría estar en riesgo y explique brevemente:
    \begin{enumerate}[label=\alph*)]
        \item Todos los individuos de la muestra tienen exactamente el mismo nivel de educación.
        \item La muestra proviene solo de trabajadores de una empresa de tecnología, pero se quiere inferir sobre todo el país.
        \item La habilidad no observada afecta salarios y está correlacionada con educación.
        \item La variabilidad del error es mayor para personas con más educación.
    \end{enumerate}

    \item \textbf{Cambio de unidades.}
    
    Suponga que se estima:
    \[
        \widehat{salario}=1000+250\,educ,
    \]
    donde $salario$ está medido en pesos por hora y $educ$ en años.
    \begin{enumerate}[label=\alph*)]
        \item Si salario se mide ahora en miles de pesos por hora, ¿cuáles son el nuevo intercepto y la nueva pendiente?
        \item Si educación se mide en semestres en vez de años, ¿cuál es la nueva pendiente?
        \item ¿Cambia el $R^2$ al cambiar las unidades? Explique.
    \end{enumerate}

    \item \textbf{Método de momentos.}
    
    En el modelo $y=\beta_0+\beta_1x+u$, suponga que se cumplen:
    \[
        E(u)=0,\qquad E(xu)=0.
    \]
    \begin{enumerate}[label=\alph*)]
        \item Escriba las contrapartes muestrales de estas dos condiciones.
        \item Muestre que estas condiciones coinciden con las condiciones de primer orden de MCO.
        \item Explique en palabras la intuición de imponer que $x$ sea ortogonal al residuo.
    \end{enumerate}

    \item \textbf{Insesgamiento y homocedasticidad.}
    
    Responda brevemente:
    \begin{enumerate}[label=\alph*)]
        \item ¿Qué significa que $\hat{\beta}_1$ sea insesgado?
        \item ¿Por qué el insesgamiento es una propiedad de la distribución muestral del estimador y no de una estimación particular?
        \item ¿Qué significa homocedasticidad?
        \item Según la fórmula
        \[
            Var(\hat{\beta}_1)=\frac{\sigma^2}{\sum_{i=1}^n(x_i-\bar{x})^2},
        \]
        ¿qué ocurre con la varianza de $\hat{\beta}_1$ si aumenta la variabilidad de $x$?
    \end{enumerate}

\end{enumerate}

\end{document}
